\section{A completion that enlarges the group}
\label{sec:hexagon}
Number the vertices of a hexagon by $i\in\mathbb Z/6\mathbb Z$, color them
by parity, and add a third-colored tip on each edge $\{i,i+1\}$ with $i$
even. This is \cite[Example 5.1]{DMV}. The three tips have degree two and
the six hexagon vertices degree three, so every correlation recognizes the
hexagon. Its allowed dihedral symmetries are exactly the stabilizer of the
alternating edge set. They are generated by $i\mapsto i+2$ and
$i\mapsto1-i$, and form $S_3$. The first generator preserves the two old
colors and the second exchanges them; the pair is $(S_3,C_3)$.

Now add a third-colored tip on all six edges. Every edge lies in its displayed
triangle, and those triangles are the maximal cliques; the completed system
is an incidence geometry. Old vertices now have degree four and tips degree
two, so again the hexagon is intrinsic. Every dihedral symmetry extends
uniquely to its incident tips and permutes the parity classes. Consequently
the full correlation group is the dihedral group of order twelve. The
color-preserving subgroup consists of rotations by even increments and
reflections $i\mapsto2a-i$. It has order six and is isomorphic to $S_3$.
In particular
\begin{equation}\label{eq:reflection}
 i\longmapsto-i,\qquad
 \text{tip on }\{i,i+1\}\longmapsto
 \text{tip on }\{-i-1,-i\}
\end{equation}
is a nonidentity type-preserving involution.

Thus the completed pair is $(D_{12},S_3)$, with the subscript denoting group
order. The source prints $(D_{12},C_3)$. Formula~\eqref{eq:reflection}
alone already excludes the second entry $C_3$; the intrinsic hexagon argument
computes the full pair. The intended obstruction survives: naive completion
of this system acquires additional correlations. The initial system already
has triangles, so it is outside the common chamber theorem. This example
does not prove that each of that theorem's hypotheses is necessary in every
possible alternative construction.
